\title{Direct Preference Optimization: Your Language Model is Secretly a Reward Model}

\begin{abstract}
Large-scale unsupervised language models (LMs) possess extensive general knowledge and exhibit some capacity for reasoning; nevertheless, guiding their output with high precision proves challenging because their training is entirely unsupervised. Current approaches to improving such controllability usually involve collecting human annotations that rank the quality of model outputs, subsequently adapting the unsupervised LM through fine-tuning to mirror these preferences, typically via reinforcement learning from human feedback (RLHF). Yet, RLHF is both intricate and prone to instability, requiring an initial phase to learn a reward model encapsulating human judgments, followed by reinforcement learning-based fine-tuning of the LM to maximize the inferred reward while minimizing deviation from the pretrained parameters. This work proposes a novel parameterization for the RLHF reward model, which allows direct derivation of the optimal policy in closed form, reducing the RLHF problem to a straightforward classification objective. This new procedure, which we refer to as \textit{Direct Preference Optimization} (DPO), is robust, efficient, and computationally lightweight. It obviates the need for sampling from the LM during fine-tuning and eliminates extensive hyperparameter searches. Empirical results indicate that DPO can fine-tune language models to conform to human preferences with performance that meets or surpasses existing techniques. Specifically, DPO demonstrates greater efficacy than PPO-based RLHF for sentiment control and matches or improves upon performance in summarization and single-turn dialogue tasks, all while offering a much simpler training and implementation process.
\end{abstract}

\section{Introduction}
Extensive unsupervised language models (LMs), trained on massive corpora, have demonstrated a range of unexpected capabilities~\citep{chowdhery2022palm, brown2020language, touvron2023llama, bubeck2023sparks}. Despite this, such models are typically exposed to data generated by humans possessing diverse intentions, expertise, and priorities. Not all of these human-generated objectives or competencies are desirable for an LM to mimic. For instance, we may wish an AI assistant specialized in coding to \textit{recognize} prevalent programming errors so it can correct them, but when it comes to code synthesis, we prefer the model to favor the (possibly infrequent) high-caliber programming skills embedded in its training corpus. Analogously, although it is valuable for a language model to \textit{identify} a misconception held by half the population, we would not want the model to propagate that misconception in half its answers when asked about it. Thus, ensuring that the model's \emph{output and conduct} are aligned with preferred outcomes—distinct from its comprehensive \textit{knowledge and potential}—is fundamental to building safe, effective, and manageable AI systems~\citep{ouyang2022training}. Although the standard approach involves aligning LMs with human preferences through reinforcement learning (RL), we will demonstrate that the RL-based objective employed by such methods can be exactly optimized via a straightforward binary cross-entropy loss, thereby streamlining the preference learning process.

Broadly, predominant approaches infuse LMs with desirable attributes by using curated datasets reflecting human preferences for beneficial and secure behaviors. This stage of preference alignment follows an initial unsupervised pre-training phase over expansive text corpora. While supervised fine-tuning on exemplary human-generated outputs provides the most direct mechanism for preference learning, reinforcement learning from human (or artificial) feedback (RLHF/RLAIF; \citep{christiano2017deep,bai2022constitutional}) has become the prevailing framework. In RLHF, a reward model is trained on annotated human preferences and subsequently used in conjunction with RL to refine an LM policy to generate outputs with high reward scores, while constraining divergence from the original pre-trained model. Although RLHF delivers language models excelling in dialogue and programming tasks, its process is substantially more intricate than supervised learning. RLHF necessitates the training of several LMs and sampling from policy distributions throughout optimization, resulting in considerable computational expense.

In this work, we demonstrate a method for aligning language model outputs with human preferences directly, sidestepping both explicit reward modeling and reinforcement learning procedures. We introduce \textit{Direct Preference Optimization (DPO)}, a novel algorithm that achieves the same optimization goal as traditional RLHF techniques—namely, reward maximization under a KL-divergence penalty—but does so through a conceptually and practically streamlined approach. The DPO algorithm modifies the model by amplifying the log-probability ratio between preferred and non-preferred answers, and introduces an adaptive, per-instance weighting factor to mitigate model collapse commonly observed with a naive log-probability ratio loss. DPO’s foundation, like that of prior algorithms, is rooted in a theoretical preference framework (for example, the Bradley-Terry model; \cite{bradley1952rankanalysis}) to evaluate the degree of correspondence between reward functions and observed preference data. Rather than using this framework to construct a preference loss for reward model training and subsequently updating the policy to optimize this learned reward, DPO leverages a variable substitution so that the preference loss is defined in terms of the policy itself. This formulation enables direct policy optimization via a standard binary cross-entropy loss, yielding a policy that is optimal under a reward function implicitly fitted to preference annotations.

Our principal contribution is the introduction of Direct Preference Optimization (DPO), a straightforward, RL-free framework for fine-tuning language models from human feedback. Experimental results demonstrate that DPO matches or surpasses the effectiveness of state-of-the-art approaches, such as PPO-based RLHF, in preference-driven applications including sentiment adjustment, text summarization, and conversational modeling, across models containing up to 6B parameters.

\section{Related Work}

As self-supervised language models scale up, they acquire the capability to handle certain tasks in a zero-shot setting \citep{radford2019language} or with minimal prompting \citep{gpt3,megatron,chowdhery2022palm}. Nonetheless, their effectiveness on practical downstream applications and their alignment with user intentions are markedly enhanced when these models undergo fine-tuning on datasets composed of explicit instructions and human-authored responses \citep{mishra-etal-2022-cross,sanh2022multitask,chung2022scaling,thoppilan2022lamda}. This approach, often termed `instruction-tuning', facilitates the generalization of LLMs to new instructions beyond those present in the fine-tuning corpus, thereby improving overall model usability \citep{chung2022scaling}. Although instruction tuning has been influential, it is often more practical to acquire \textit{relative} human feedback regarding the quality of responses rather than expert-labeled examples. Motivated by this, recent studies have introduced fine-tuning protocols that employ datasets derived from human preferences, yielding improvements in tasks such as translation \citep{kreutzer-etal-2018-reliability}, summarization \citep{stiennon2022learning,ziegler2020finetuning}, narrative generation \citep{ziegler2020finetuning}, and adherence to instructions \citep{ouyang2022training,ramamurthy2023is}. In these methods, a neural network is first trained as a reward function to reflect the collected preference data, using frameworks like the Bradley-Terry model \citep{bradley1952rankanalysis}. The language model is subsequently refined to maximize this learned reward function, typically employing reinforcement learning methods such as REINFORCE \citep{williams1992reinforce}, proximal policy optimization (PPO; \cite{schulman2017proximal}), or related algorithms \citep{ramamurthy2023is}. A related area of research exploits LLMs that have been instruction-tuned with human-in-the-loop feedback to generate synthetic preference annotations focusing on particular characteristics, for example, safety or harmlessness \citep{bai2022constitutional}, utilizing light-touch human supervision through guideline rubrics for the model's annotations. These developments unify two distinct streams of research: the application of reinforcement learning to train language models across diverse objectives~\citep{Ranzato2015SequenceLT,paulus2018a,wu2018learning}, and general strategies for learning from human preference data \citep{christiano2017deep,kupcsik2018learning}. However, despite the advantages of leveraging comparative human judgments, practical obstacles persist in fine-tuning large language models with reinforcement learning. This work addresses these challenges by presenting a theoretically grounded method to optimize relative human preferences without reliance on reinforcement learning.

Research on learning policies from preferences in non-linguistic domains has been conducted in both bandit and reinforcement learning frameworks, yielding a variety of methods. Within contextual bandit settings, the study of learning from preferences or rankings, as opposed to explicit reward signals, is known as the contextual dueling bandit (CDB) problem \citep{yue2012karmed, dudik2015contextual}. When explicit reward information is unavailable, theoretical work on CDBs replaces the concept of an optimal policy with the \textit{von Neumann winner}, which refers to a policy that achieves at least a 50\% expected win rate against all possible alternative policies \citep{dudik2015contextual}. Notably, in the CDB paradigm, preferences are typically provided in an online fashion, whereas the standard setup in learning from human preferences involves leveraging a static offline dataset of action pairs annotated with preferences \citep{yan2022human}. Analogously, \textit{preference-based RL} (PbRL) focuses on learning from binary preference information generated by an implicit, unobserved ‘scoring’ mechanism rather than explicit rewards \citep{BusaFekete2014, ruiz2023dueling}. Existing PbRL techniques, including those capable of utilizing off-policy preference data, predominantly proceed in two stages: they first infer the underlying scoring function (that is, build a reward model), followed by its optimization \citep{jain2013learning, BusaFekete2014, christiano2017deep, sadigh2017active, kupcsik2018learning}. Contrasting with this prevailing approach, we propose a unified policy learning framework that optimizes directly for preference satisfaction without separate reward model estimation.

\section{Preliminaries}\label{section:prelims}

We summarize the RLHF pipeline introduced by \citeauthor{ziegler2020finetuning} and further developed in subsequent work \citep{stiennon2022learning, bai2022training, ouyang2022training}. This process is typically partitioned into three sequential stages: (1) supervised fine-tuning (SFT); (2) sampling preferences and training a reward model; and (3) reinforcement learning-based optimization.

\textbf{SFT}: The RLHF process typically initiates with the supervised fine-tuning of a pre-trained language model on curated datasets tailored to the target application (e.g., dialogue, summarization), resulting in an updated model denoted $\pi^\text{SFT}$.

\textbf{Reward Modelling Phase}: In the subsequent phase, the fine-tuned SFT model generates response pairs $(y_1, y_2)\sim \pi^\text{SFT}(y \mid x)$ for given prompts $x$. Human annotators are then tasked with indicating a preference between each pair, yielding annotations of the form $y_w\succ y_l \mid x$, where $y_w$ and $y_l$ denote the preferred and non-preferred completions, respectively, for each prompt. These human preferences are modeled as being governed by a latent, unknown reward function $r^*(y, x)$. To characterize the process underlying these preferences, several probabilistic models have been considered, with the Bradley-Terry (BT) model \cite{bradley1952rankanalysis} being widely used. More general Plackett-Luce models \citep{plackett1975analysis, luce2012individual} are also applicable when multiple ranked outputs are available. The BT model formalizes the human preference distribution $p^*$ as:

\begin{equation}\label{eq:bradley-terry}
    p^*(y_1\succ y_2 \mid x)=\frac{\exp\left(r^*(x, y_1)\right)}{\exp\left(r^*(x, y_1)\right) + \exp\left(r^*(x, y_2)\right)}.
\end{equation}

Given a fixed dataset of comparison samples $\mathcal{D}=\bigl\{x^{(i)}, y_w^{(i)}, y_l^{(i)}\bigr\}_{i=1}^N$ drawn according to $p^*$, one can represent the reward model $r_{\phi}(x, y)$ and learn its parameters by maximizing likelihood. Casting this as a binary classification task, the associated negative log-likelihood objective becomes:

\begin{equation}\label{eq:reward_model}
    \mathcal{L}_R(r_{\phi}, \mathcal{D}) = -\mathbb{E}_{(x, y_w, y_l)\sim \mathcal{D}}\bigl[\log \sigma(r_{\phi}(x, y_w)- r_{\phi}(x, y_l))\bigr]
\end{equation}

with $\sigma$ denoting the logistic sigmoid. For language modeling applications, $r_{\phi}(x, y)$ is usually initialized from the SFT policy $\pi^\text{SFT}(y \mid x)$, augmented by a linear output head after the transformer, producing a single scalar reward value \cite{ziegler2020finetuning}. To limit variance in reward outputs, prior approaches standardize the rewards so that $\mathbb{E}_{x,y\sim \mathcal{D}}\left[r_\phi(x, y)\right] = 0$ across all $x$.

\textbf{RL Fine-Tuning Phase}: During RL-based fine-tuning, this learned reward provides supervision for the language model. Consistent with earlier works~\citep{jaques2017sequence, jaques2020human}, the policy is optimized according to

\begin{equation}\label{eq:RL}
\max_{\pi_{\theta}}  \mathbb{E}_{x\sim \mathcal{D}, y\sim \pi_{\theta}(y \mid x)}\bigl[r_{\phi}(x, y)\bigr] - \beta\mathbb{D}_{\textrm{KL}}\bigl[\pi_{\theta}(y\mid x)\mid \mid \pi_\text{ref}(y\mid x)\bigr],
\end{equation}

where the coefficient $\beta$ mediates how much the optimized policy $\pi_\theta$ is regularized to stay close to the reference policy $\pi_\text{ref}$, typically taken as the initial SFT model $\pi^\text{SFT}$. Practically, $\pi_\theta$ is also initialized from $\pi^\text{SFT}$. This KL regularization is crucial, as it ensures that the policy does not stray too far from regions where the reward model predictions are reliable, thus safeguarding diversity in the model's outputs and mitigating the risk of mode-collapse onto only the highest-reward sequences. Since text generation operates in a discrete space, direct optimization is not differentiable, and reinforcement learning techniques are employed. The conventional procedure \citep{ziegler2020finetuning, stiennon2022learning, bai2022training, ouyang2022training} defines the reward as ${r(x, y) = r_{\phi}(x, y) -\beta (\log \pi_{\theta}(y\mid x) - \log \pi_\text{ref}(y\mid x))}$, with the policy then improved using PPO \cite{schulman2017proximal}.

\section{Direct Preference Optimization}\label{sec:DPO}

Recognizing the limitations of reinforcement learning when scaling to large problems such as language model fine-tuning, we aim to develop a straightforward preference-based policy optimization method. Distinct from conventional RLHF methods that separately learn a reward function before optimizing it via RL, our technique exploits a specific reward model structure that admits a closed-form optimal policy, bypassing the need for an RL training loop. The core insight is an analytic mapping from reward functions to their corresponding optimal policies, which permits the substitution of a reward-oriented objective for an equivalent policy-oriented one. Through this transformation, our approach sidesteps the need to fit an explicit reward model, while still adhering to human preference structures, as characterized by models such as Bradley-Terry. In this way, the policy serves both as the language model and implicitly encodes the reward.

\textbf{DPO Objective Derivation.} To begin, we adopt the same reinforcement learning objective as prior literature, Eq.~\ref{eq:RL}, with a generic reward function $r$. As demonstrated in previous studies~\citep{peters2007reinforcement, peng2019advantage, korbak2022reinforcement, go2023aligning}, the solution that optimizes the KL-regularized reward maximization problem described in Eq.~\ref{eq:RL} can be written as:

\begin{equation}\label{eq:op_policy}
    \pi_r(y\mid x) = \frac{1}{Z(x)}\pi_\text{ref}(y\mid x)\exp\left(\frac{1}{\beta}r(x, y)\right),
\end{equation}

where $Z(x) = \sum_{y}\pi_\text{ref}(y\mid x)\exp\left(\frac{1}{\beta}r(x, y)\right)$ serves as the normalization constant or partition function. Detailed derivations can be found in Appendix \ref{app:derivation1}. Notably, estimating $Z(x)$ remains computationally intensive—even when employing an MLE-based reward estimator $r_{\phi}$ as a stand-in for the true reward $r^*$~\citep{korbak2022reinforcement, go2023aligning}—posing practical challenges for this formulation. Nonetheless, Eq.~\ref{eq:op_policy} can be rearranged so that the reward function is re-expressed via the optimal policy $\pi_r$, the reference distribution $\pi_\text{ref}$, and the as-yet-unknown $Z(\cdot)$. By applying the natural logarithm to both sides of Eq.~\ref{eq:op_policy}, we derive the following expression after simplification:

\begin{equation}\label{eq:main_eq}
    r(x,y) = \beta \log \frac{\pi_r(y\mid x)}{\pi_\text{ref}(y\mid x)} + \beta \log Z(x).
\end{equation}

This transformation can be equivalently applied to the actual reward function $r^*$ and its optimal policy $\pi^*$. Importantly, the Bradley-Terry model is sensitive only to the difference between rewards of two choices, specifically, ${p^*(y_1 \succ y_2 \mid x) = \sigma(r^*(x, y_1) - r^*(x, y_2))}$. Replacing $r^*(x, y)$ with its expression from Eq.~\ref{eq:main_eq} in the preference model of Eq.~\ref{eq:bradley-terry}, the dependence on the partition function $Z(x)$ vanishes. As a result, the human preference probability can be articulated solely in terms of the optimal policy $\pi^*$ and the reference policy $\pi_\text{ref}$. Consequently, for the Bradley-Terry model, the optimal RLHF policy $\pi^*$ conforms to:

\begin{equation}\label{eq:objective}
    p^*(y_1\succ y_2 \mid x) = \frac{1}{1 + \exp\left(\beta \log \frac{\pi^*(y_2\mid x)}{\pi_\text{ref}(y_2\mid x)} - \beta \log \frac{\pi^*(y_1\mid x)}{\pi_\text{ref}(y_1\mid x)}\right)}
\end{equation}

Details of this derivation are provided in Appendix~\ref{app:derivation2}. While the above relies on the Bradley-Terry framework, an analogous development can be carried out for the Plackett-Luce models~\citep{plackett1975analysis, luce2012individual}, with explicit forms detailed in Appendix~\ref{app:plackett_luce_models}.

With this probabilistic representation of preference grounded in the optimal policy, rather than the reward, it becomes feasible to optimize a parameterized policy $\pi_\theta$ using a maximum likelihood principle. Drawing a parallel to reward modeling (see Eq.~\ref{eq:reward_model}), the policy learning objective is reformulated as follows:

\begin{equation}\label{eq:optimum_model}
    \mathcal{L}_\text{DPO}(\pi_{\theta}; \pi_\text{ref}) = -\mathbb{E}_{(x, y_w, y_l)\sim \mathcal{D}}\left[\log \sigma \left(\beta \log \frac{\pi_{\theta}(y_w\mid x)}{\pi_\text{

In this approach, we estimate an implicit reward through an alternative parameterization, where the optimal policy directly corresponds to $\pi_\theta$. Additionally, since our method is mathematically analogous to training a reparameterized Bradley-Terry model, it inherits several theoretical guarantees—including consistency under certain conditions on the distribution of preference data, as discussed by \cite{bong2022generalized}. We expand on the theoretical underpinnings of DPO and its connection to existing literature in Section~\ref{sec:theory}.

\textbf{How does the DPO update operate?} To gain mechanistic insight into DPO, we can inspect the gradient of its loss function, $\mathcal{L}_\text{DPO}$. The derivative of the loss with respect to the model parameters $\theta$ can be expressed as:

\begin{multline*}\label{eq:gradient}
    \nabla_\theta \mathcal{L}_\text{DPO}(\pi_\theta;\pi_\text{ref}) = \\ -\beta\mathbb{E}_{(x, y_w, y_l) \sim \mathcal{D}} \bigg[\underbrace{\sigma(\hat{r}_\theta(x, y_l) - \hat{r}_\theta (x, y_w))}_\text{larger when reward assignment is incorrect}\bigg[\underbrace{\nabla_\theta\log \pi(y_w \mid x)}_\text{promote $y_w$ occurrence} - \underbrace{\nabla_\theta\log\pi(y_l \mid x)}_\text{reduce $y_l$ occurrence}\bigg]\bigg],
\end{multline*}

Here, $\hat{r}_\theta(x, y) = \beta \log \frac{\pi_\theta(y \mid x)}{\pi_\text{ref}(y \mid x)}$ defines the implicit reward function as determined by the language model $\pi_\theta$ relative to the reference model $\pi_\text{ref}$ (additional details are provided in Section~\ref{sec:theory}). In essence, the gradient of $\mathcal{L}_\text{DPO}$ acts to boost the probability of preferred responses $y_w$ and suppress the probability of non-preferred responses $y_l$. Notably, each training example is weighted according to the degree by which the implicit reward model $\hat{r}_\theta$ mistakenly favors the dispreferred option, with this discrepancy modulated by $\beta$—thus capturing both the extent and magnitude of ordering errors, as well as the influence of the KL regularization. Our empirical findings highlight the significance of this weighting scheme: omitting the coefficient and using an unweighted approach leads to poor language model behavior (see Appendix Table~\ref{tab:unlikelihood_generations}).

\textbf{DPO overview.} The typical workflow for DPO consists of: 1) For each prompt $x$, generate completion pairs $y_1, y_2 \sim \pi_\text{ref}(\cdot \mid x)$, then annotate these with human preferences to assemble an offline preference dataset $\mathcal{D} = \{x^{(i)}, y_w^{(i)}, y_l)^{(i)}\}_{i=1}^N$; 2) train the language model $\pi_\theta$ by minimizing the loss $\mathcal{L}_\text{DPO}$, using the specified $\pi_\text{ref}$, dataset $\mathcal{D}$, and tuning parameter $\beta$.  
In real-world scenarios, researchers generally prefer to leverage existing, publicly available preference datasets instead of producing new completions and soliciting human feedback. Because these datasets are typically generated with $\pi^\text{SFT}$, the reference policy is set as $\pi_\text{ref} = \pi^\text{SFT}$ whenever it is obtainable. If $\pi^\text{SFT}$ is unavailable, the reference model $\pi_\text{ref}$ is instead initialized by maximizing the likelihood over the preferred completions ${(x, y_w)}$ according to ${\pi_\text{ref} = \argmax_{\pi}\mathbb{E}_{x, y_w \sim \mathcal{D}}\left[\log \pi(y_w \mid x)\right]}$. This approach reduces the discrepancy between the unknown ideal reference distribution and the practically used $\pi_\text{ref}$ in DPO. Further information about implementation specifics and hyperparameter selection is provided in Appendix~\ref{app:implementation}.

\section{Theoretical Analysis of DPO}  
This section presents a deeper theoretical understanding of the DPO approach, along with an explanation of its properties and a comparison highlighting its benefits relative to actor-critic methods in RLHF—such as PPO~\cite{schulman2017proximal}—which face their own challenges.

\label{sec:theory}

\subsection{Your Language Model Is Secretly a Reward Model}  
DPO achieves preference optimization without the need for explicit reward function learning or separate RL policy optimization steps, by relying solely on a maximum likelihood training regime. Notably, the objective from Eq. \ref{eq:main_eq} corresponds to the Bradley-Terry model with a reward specified as $r^*(x, y) = \beta \log\frac{\pi^*_\theta(y \mid x)}{\pi_\text{ref}(y \mid x)}$; thus, optimizing the parametric language model $\pi_{\theta}$ can be viewed as analogous to optimizing the reward model as in Eq. \ref{eq:reward_model}, after a suitable change of variables. The subsequent analysis will formalize this reparameterization, demonstrate that it preserves the full expressivity for the class of reward models being learned, and prove that it enables precise recovery of the optimal policy. To proceed, we first establish an equivalence relation among reward functions.

\begin{definition}
Two reward functions $r(x, y)$ and $r'(x, y)$ are defined as equivalent if and only if there exists a function $f$ such that ${r(x, y)-r'(x, y) = f(x)}$.
\end{definition}

It is straightforward to verify that this definition constitutes an equivalence relation, resulting in a partition of the set of all reward functions into distinct equivalence classes. The following lemmas can then be formulated:

\begin{lemma}\label{lemma:same_prefrence}
Within the Plackett-Luce framework, including the Bradley-Terry preference model as a specific instance, any two reward functions belonging to the same equivalence class lead to an identical distribution over preferences.
\end{lemma}

\begin{lemma}\label{lemma:same_policy}
For the constrained reinforcement learning setting, all reward functions from the same equivalence class yield the same optimal policy.
\end{lemma}

These statements follow directly and their proofs are postponed to Appendix \ref{app:lemma1}. The first lemma highlights the well-known identifiability challenge in Plackett-Luce models \cite{plackett1975analysis}: without further restrictions, the parameterization is not unique. Consequently, it is standard practice to impose identifiability constraints to ensure the maximum likelihood estimates derived from Eq. \ref{eq:reward_model} are meaningful \cite{bong2022generalized}. The second lemma asserts that all members of a given equivalence class produce the same optimal solution, implying that the learning goal reduces to finding any representative of the equivalence class that aligns with the optimal policy. We formalize this principle in the following theorem, proved in Appendix~\ref{app:thm1}:

\begin{theorem}\label{thm:main}
Assuming mild technical conditions, every reward equivalence class compatible with the Plackett-Luce (including the Bradley-Terry model) can be expressed by the parameterization ${r(x, y) = \beta \log \frac{\pi(y\mid x)}{\pi_\text{ref}(y\mid x)}}$ for some model $\pi(y\mid x)$ with respect to a fixed reference model $\pi_\text{ref}(y \mid x)$.
\end{theorem}

\begin{sproof}
Take any reward function $r(x, y)$ that determines an associated optimal model $\pi_r(y \mid x)$ as described by Eq. \ref{eq:op_policy}. We aim to demonstrate that within the equivalence class of $r$, there is a reward function that can be reparameterized as stated. Let us define the mapping $f$:
\begin{equation}
    f(r; \pi_\text{ref}, \beta)(x, y) = r(x, y) - \beta\log\sum_{y}\pi_\text{ref}(y\mid x)\exp\left(\frac{1}{\beta}r(x, y)\right)
\end{equation}
This operator $f$ renormalizes the original reward by subtracting the log-partition function associated with $\pi_r$, depending solely on $x$. Since this adjustment introduces only a function of $x$, $f(r; \pi_\text{ref}, \beta)(x, y)$ remains in the equivalence class of $r(x, y)$. Substituting $r$ using the right side of Eq.~\ref{eq:main_eq} (valid for any reward), it follows that $f(r; \pi_\text{ref}, \beta)(x, y) = \beta \log \frac{\pi_r(y\mid x)}{\pi_\text{ref}(y\mid x)}$. Therefore, $f$ yields a canonical representative of $r$’s equivalence class, demonstrating that the reparameterization is without loss of generality.
\end{sproof}

An alternative perspective on Theorem~\ref{thm:main} is that it precisely identifies the particular reward function, among all those within a given equivalence class, that is selected by the DPO reparameterization. Specifically, it picks out the reward function that fulfills the following condition:

\begin{equation}\label{eq:lag_p}
     \sum_{y}\underbrace{\pi_\text{ref}(y\mid x)\exp\left(\frac{1}{\beta}r(x, y)\right)}_{=\pi(y\mid x)\text{, using Thm.~\ref{thm:main} reparam.}} = 1,
\end{equation}

which ensures that $\pi(y\mid x)$ constitutes a legitimate probability distribution (i.e., all probabilities are non-negative and sum to one). By comparing with Eq.~\ref{eq:op_policy}, it is evident that Eq.~\ref{eq:lag_p} defines the partition function corresponding to the optimal policy that emerges from the reward $r(x, y)$. The DPO algorithm’s crucial insight lies in the possibility of applying suitable constraints to the otherwise under-determined Plackett-Luce (and particularly Bradley-Terry) class of preference models. This allows the set of representable reward models to be maintained while also ensuring that the optimal policy given in Eq. \ref{eq:op_policy} is analytically manageable across all prompts $x$.

\subsection{Instability of Actor-Critic Algorithms}

Our framework is also effective for analyzing sources of instability in widely-used actor-critic algorithms for RLHF, such as PPO. Focusing on the RL fine-tuning process detailed in Section \ref{section:prelims}, we connect our analysis to the control as inference paradigm \cite{levine2018reinforcement} as applied to the constrained RL problem formulated in \ref{eq:RL}. Letting $\pi_{\theta}(y\mid x)$ be a parametric policy, the approach aims to minimize $\mathbb{D}_{\text{KL}}[\pi_{\theta}(y|x) \mid \mid \pi^*(y\mid x)]$, where $\pi^*$, defined in Eq. \ref{eq:optimum_model}, denotes the optimal policy corresponding to the reward function $r_{\phi}(y, x)$. Performing the necessary derivations, the resulting optimization problem can be expressed as:

\begin{equation}\label{eq:AC}
    \max_{\pi_{\theta}}\mathbb{E}_{\pi_{\theta}(y\mid x)}\bigg[\underbrace{r_{\phi}(x, y) -\beta\log\sum_{y}\pi_\text{ref}(y\mid x)\exp\left(\frac{1}{\beta}r_{\phi}(x, y)\right)}_{f(r_{\phi}, \pi_\text{ref}, \beta)} - \underbrace{\beta\log\frac{\pi_{\theta}(y\mid x)}{\pi_\text{ref}(y\mid x)}}_{\text{KL}}\bigg]
\end{equation}

The objective under consideration aligns with that optimized in earlier studies 
\citep{ziegler2020finetuning, stiennon2022learning, bai2022training, ouyang2022training}, making use of the DPO-equivalent reward within the reward class $r_{\phi}$. Within this context, the normalization component in $f(r_{\phi}, \pi_\text{ref}, \beta)$ can be interpreted as the soft value function associated with the reference policy $\pi_\text{ref}$. Although omitting this normalization does not influence the optimality of the solution, its absence can induce substantial variance in policy gradient estimates, leading to potential instability during training. One strategy to handle this normalization term involves utilizing a learned value function; however, this approach often presents its own optimization challenges. Prior methods have addressed normalization by referencing a human completion baseline, which effectively serves as a Monte Carlo estimate based on a single sample for the normalization term. Unlike these strategies, the DPO reparameterization inherently constructs a reward function that is independent of explicit baselines.

\section{Experiments}
This section presents an empirical investigation into the effectiveness of DPO for learning policies directly from preference data. We first analyze DPO in a controlled text generation environment, specifically probing how it balances reward maximization with KL-divergence minimization relative to established preference learning methods such as PPO. Subsequently, we assess DPO on more challenging RLHF scenarios and with larger models, including applications in summarization and conversational tasks. Across these settings, we observe that DPO, even with minimal hyperparameter adjustment, typically matches or surpasses the performance of strong benchmarks, including RLHF with PPO and the strategy of selecting the best out of $N$ trajectories under a learned reward model. Prior to detailing our findings, we outline the experimental protocols; further specifics can be found in Appendix~\ref{app:exp_details}.

\textbf{Tasks.} Our study investigates three distinct open-ended text generation scenarios. In all cases, algorithms acquire a policy using a dataset of preferences $\mathcal{D}=\bigl\{x^{(i)}, y_w^{(i)}, y_l^{(i)}\bigr\}_{i=1}^N$. The \textbf{controlled sentiment generation} task employs $x$ as a prompt extracted from a movie review in the IMDb dataset \cite{maas-EtAl:2011:ACL-HLT2011}, and the objective for the policy is to produce $y$ exhibiting positive sentiment. For systematic evaluation, preference pairs over generated outputs are synthesized using a pre-trained sentiment classifier such that $p(\text{positive}\mid x,y_w)>p(\text{positive}\mid x,y_l)$. For SFT, GPT-2-large is fine-tuned on training set reviews from IMDB until convergence (additional methodological specifics appear in App~\ref{app:sentiment_details}). The \textbf{summarization} task involves $x$ as a Reddit forum post, with the policy tasked to produce a summary $y$ that captures the main ideas. Consistent with existing literature, we utilize the Reddit TL;DR summarization dataset \citep{volske-etal-2017-tl} in combination with human preference data collected by \citeauthor{stiennon2022learning}. For RLHF, we employ an SFT model already fine-tuned on human-authored summaries of forum posts\footnote{\url{https://huggingface.co/CarperAI/openai_summarize_tldr_sft}} through the TRLX \citep{leandro_von_werra_2023_7790115} toolkit. Note that the preference labels were generated by \citeauthor{stiennon2022learning} using outputs from a different but comparably trained SFT system. The \textbf{single-turn dialogue} task frames $x$ as a user inquiry, ranging from scientific to personal domains. Here, the policy must craft a helpful and engaging reply $y$; we utilize the Anthropic Helpful and Harmless dialogue dataset \citep{bai2022training}, which includes 170k conversational exchanges between a person and an automated agent. In every case, two responses are presented at the conclusion of each transcript, generated by a sizable language model and annotated with a human preference indicator for the superior answer. In this task, there is no readily available SFT model; we therefore construct the SFT model by fine-tuning a general-purpose language model solely on human-preferred completions.

\textbf{Evaluation.} We utilize two distinct evaluation methodologies in our experiments. To scrutinize how well each algorithm optimizes the constrained reward maximization objective, we examine their respective frontiers of achieved reward and KL-divergence from the reference policy under the controlled sentiment generation scenario. This analysis is feasible due to access to a ground-truth reward function, instantiated as a sentiment classifier. In contrast, real-world applications lack such oracle reward functions. Consequently, we assess algorithms by their \textit{win rate} against a baseline policy, with GPT-4 acting as a surrogate for human assessment in both summary quality and response helpfulness within the summarization and single-turn dialogue tasks, respectively. For summarization, we take the test set's reference summaries as the comparator baseline, whereas for dialogue, the dataset’s annotated preferred responses serve this role. Although large language models have demonstrated stronger correspondence to human judgments than traditional automated metrics \citep{Chen2023ExploringTU}, we further support our choice of GPT-4 for evaluation with a dedicated human subject study (see Sec.~\ref{sec:human-judgments}). Our results indicate that GPT-4's evaluations display high concordance with human raters, often matching or exceeding the level of agreement observed among human annotators themselves.

\textbf{Methods.} Beyond DPO, our experimental suite includes several established methods for training language models to align with human preference data. We examine zero-shot prompting on the summarization task with \textbf{GPT-J} \citep{gpt-j}, and two-shot prompting for dialogue using \textbf{Pythia-2.8B} \citep{biderman2023pythia}. We also assess the \textbf{SFT} model, as well as \textbf{Preferred-FT}, the latter being a supervised model trained on the selected completion $y_w$, derived from SFT outputs (for sentiment and summarization) or a generic language model (for dialogue). Additionally, we include \textbf{Unlikelihood}~\citep{welleck2019neural}, which seeks to increase the likelihood of $y_w$ while explicitly reducing the likelihood of $y_l$, utilizing an optional weighting factor $\alpha \in [0, 1]$ on the “unlikelihood” component. Our comparisons also extend to \textbf{PPO} \citep{schulman2017proximal}, trained with a reward model inferred from preference data, and \textbf{PPO-GT}, an oracle leveraging the exact reward function in the controlled sentiment setup. In sentiment-based trials, we adopt two PPO-GT implementations: an off-the-shelf version \cite{leandro_von_werra_2023_7790115}, and an adapted variant that incorporates reward normalization and hyperparameter tuning for enhanced results; these modifications are also applied in PPO runs with learned rewards. Finally, we employ a \textbf{Best of $N$} approach: from either the SFT model (or Preferred-FT in dialogue), $N$ completions are generated, and the one scoring highest according to the learned reward function is selected. Although this technique can yield strong performance, its dependency on evaluating $N$ candidate completions per query during inference makes it computationally unscalable as $N$ grows.

\subsection{How well can DPO optimize the RLHF objective?}

Typical RLHF methods rely on a KL-constrained reward maximization framework that seeks to optimize the reward function while ensuring that the learned policy does not diverge excessively from the reference policy. Consequently, comparisons between methods should consider both the attained reward and the KL divergence; substantially increasing reward at the expense of a large KL is often counterproductive. Figure~\ref{fig:frontier-tldr-main} illustrates the trade-off between reward and KL for different algorithms under the sentiment task. For each method, we conduct several training runs, varying the policy conservativeness hyperparameter (for PPO, target KL $\in\{3,6,9,12\}$; for DPO, $\beta \in \{0.05,0.1,1,5\}$; for unlikelihood, $\alpha\in\{0.05,0.1,0.5,1\}$; and for preferred-FT, different random seeds), resulting in 22 training runs in total. After every 100 steps throughout training, policies are evaluated on a fixed set of test prompts, calculating the mean reward based on the ground truth reward function along with the mean sequence-level KL\footnote{This refers to the cumulative KL-divergence computed over all timesteps in the sequence.} relative to the reference policy $\text{KL}\left(\pi\mid \mid \pi_\text{ref}\right)$. Our findings demonstrate that DPO establishes the most favorable reward-KL frontier, attaining superior reward with correspondingly low KL values. This outcome is noteworthy for several reasons. Firstly, both DPO and PPO target identical optimization objectives, yet DPO consistently demonstrates greater efficiency; the reward/KL frontier for DPO dominates that of PPO. Secondly, DPO outperforms PPO not only in general, but also surpasses PPO even when PPO is supplied with access to exact ground truth rewards (PPO-GT).

\subsection{Can DPO scale to real preference datasets?}
\label{sec:dpo-real-datasets}

We proceed to assess how well DPO performs when fine-tuned on real-world preference data for summarization and single-turn dialogue tasks. In the case of summarization, commonly used automatic metrics such as ROUGE often exhibit weak correlation with actual human judgment~\citep{stiennon2022learning}, and earlier research has shown that fine-tuning language models using PPO on human feedback results in improved summary quality. Our evaluation contrasts several approaches by generating completions on the TL;DR summarization dataset’s test split and calculating the mean win rate relative to test reference completions. To gauge robustness, completions from all models are produced across a spectrum of sampling temperatures ranging from 0.0 to 1.0, with win rates displayed in Figure~\ref{fig:frontier-tldr-main} (right). The models—including DPO, PPO, and Preferred-FT—are all fine-tuned starting from the identical GPT-J SFT checkpoint\footnote{\url{https://huggingface.co/CarperAI/openai_summarize_tldr_sft}}. Experimental results indicate that DPO achieves a win rate of about 61\% at temperature 0.0, outperforming PPO, which peaks at around 57\% at the same temperature. Furthermore, DPO secures a superior maximum win rate in comparison to the best $N$ baseline. It should be highlighted that DPO's $\beta$ hyperparameter was not extensively optimized for these experiments, suggesting possible underestimation of DPO's performance ceiling. DPO also demonstrates greater resilience to changes in sampling temperature, whereas PPO's effectiveness diminishes, eventually matching the base GPT-J model at higher temperatures. Notably, Preferred-FT yields minimal gains over the initial SFT model. For a direct comparison, Section~\ref{sec:human-judgments} presents human evaluation data showing that DPO outputs at temperature 0.25 are selected over PPO outputs (also at temperature 0.25) 58\% of the time.

For single-turn dialogue, we assess various approaches on the subset of the Anthropic HH test split \citep{bai2022training} that consists of a single human-assistant exchange. For GPT-4 evaluations, we determine win rates by comparing generated outputs to the test’s preferred completions, which serve as references. In the absence of a conventional SFT baseline, we initialize with a pre-trained Pythia-2.8B model, train a reference model via Preferred-FT on the preferred completions (ensuring the resulting outputs remain in-distribution), and subsequently apply DPO training. Comparisons are drawn against both the best out of 128 Preferred-FT completions (where performance of the Best of $N$ method saturates at $N=128$ for this scenario; refer to Appendix Figure~\ref{fig:best-of-n}) and a 2-shot prompt setup using the original Pythia-2.8B, observing that DPO matches or exceeds the best outcomes for each approach across optimal temperature settings. Additionally, we evaluate an RLHF agent fine-tuned using PPO on the Anthropic HH set\footnote{\url{https://huggingface.co/reciprocate/ppo_hh_pythia-6B}}, as made available by a reputable source\footnote{\url{https://github.com/CarperAI/trlx/tree/main/examples/hh}}. However, no tested prompts or temperature configurations allow this RLHF model to outperform the baseline Pythia-2.8B. Drawing on findings from TL;DR and given that both Best of 128 and PPO approaches aim to maximize the same reward objective, we interpret the Best of 128 baseline as a reasonable stand-in for PPO-level capability. Ultimately, DPO emerges as the sole method that both efficiently and effectively surpasses the preferred completion benchmark in the Anthropic HH corpus, and it achieves comparable or superior performance relative to the much more resource-intensive Best of 128 baseline. Notably, as illustrated in Figure~\ref{fig:dialogue-main}, DPO reaches optimal performance with relatively rapid convergence.

\subsection{Generalization to a new input distribution}

To further investigate how PPO and DPO respond to shifts in input distribution, we evaluate both sets of policies—trained on the Reddit TL;DR summarization task—on an alternative domain: news articles from the test split of the CNN/DailyMail dataset \citep{nallapati-etal-2016-abstractive}. We use the optimal sampling temperatures previously established on TL;DR (0 and 0.25) and report the outcomes in Table~\ref{tab:ood}. For assessment, we calculate the win rate of GPT-4 when pitted against the dataset’s reference summaries, employing the same GPT-4 (C) prompt as in the Reddit TL;DR evaluations, with the term ``forum post'' replaced by ``news article.'' On this out-of-domain data, DPO continues to surpass the PPO policy by a large margin. These results provide preliminary support for the conclusion that DPO achieves generalization comparable to PPO, even though PPO leverages additional unlabeled Reddit TL;DR prompts that are not used by DPO.

\subsection{Validating GPT-4 judgments with human judgments}
\label{sec:human-judgments}
We implement a human evaluation to assess the credibility of GPT-4's preferences, focusing on the TL;DR summarization task and employing two different prompting strategies. The \textbf{GPT-4 (S)} (simple) prompt asks which summary best captures the crucial information from the post, whereas the \textbf{GPT-4 (C)} (concise) prompt additionally queries about summary conciseness; the latter is particularly relevant, as GPT-4 with the \textbf{GPT-4 (S)} prompt often favors longer and more repetitive summaries than those preferred by human raters. Full prompt details are provided in Appendix~\ref{app:prompts}. We perform three specific comparisons: the top-performing DPO model (temperature 0.25), the lowest-performing PPO model (temperature 1.0), and an intermediate SFT model (temperature 0.25), all of which are evaluated against the best-performing, greedily-sampled PPO variant. This range is selected to ensure coverage of various sample qualities; for each of the three cases, both prompts are tested. Our findings reveal that, regardless of prompt, GPT-4’s assessments correspond with human judgments at a rate comparable to inter-human agreement, indicating GPT-4's suitability as a stand-in for human evaluators. Due to a limited pool of raters, multiple human judgments are collected only for the DPO and PPO-1 model comparisons. On the whole, the \textbf{GPT-4 (C)} prompt yields win rates that align more closely with those of human raters, and thus we utilize this prompt for the principal experiments in Section~\ref{sec:dpo-real-datasets}. More comprehensive information about the human study—including interface design and the roster of participants—can be found in Appendix~\ref{app:human-study}.

\section{Discussion}
Preference-based learning stands out as a robust and adaptable strategy for developing capable and value-aligned language models. In this work, we have proposed DPO, a streamlined approach that allows language models to be trained from preference data without employing reinforcement learning. Instead of transforming the preference learning challenge into a conventional RL scenario to utilize existing RL algorithms, DPO establishes a direct correspondence between reward functions and language model policies. This enables direct optimization towards human preferences using a simple cross-entropy loss, bypassing both reinforcement learning frameworks and any restrictions on generality. DPO achieves performance on par with or surpassing established RLHF techniques, such as those based on PPO, requiring minimal hyperparameter adjustment. This substantially lowers the practical difficulty of leveraging human preference data to train additional language models.

\textbf{Limitations \& Future Work.} Our findings highlight numerous avenues for continued investigation. One pertinent question is the extent to which DPO-trained policies can generalize beyond the training distribution in comparison to methods relying on explicit reward learning. Preliminary experiments imply that DPO generalizes as well as models optimized with PPO, yet a thorough analysis remains necessary. Furthermore, the potential of using DPO-based self-labeling to make effective use of prompts without human labels merits deeper examination. Another aspect for future inquiry involves understanding how phenomena such as reward over-optimization appear in DPO, and whether the minor dip in results observed in Figure~\ref{fig:dialogue-main}-right is attributable to this issue. Our evaluations have primarily considered models up to 6B parameters; scaling DPO to even larger, state-of-the-art architectures presents a compelling direction for subsequent research. In terms of evaluation methodology, we noticed that the win rates as judged by GPT-4 vary depending on the prompt, suggesting future studies should investigate optimal strategies for extracting reliable evaluations from automated systems. Beyond language modeling, DPO's methodology has potential applications for training generative models across a range of domains informed by human preferences.